% English counterpart of draft/background_and_motivation.md.
% Compile from paper/; see motivation-preview.tex for EuroSys 2027 settings.
\section{Background and Motivation}
\label{sec:background-motivation}

Multi-agent applications complete a user task through interdependent LLM
requests. After a worker failure, rebuilding lost key--value (KV) state can
restore an individual request, but the user-visible penalty depends on when
the entire workflow completes. Workflow dependencies and the cost of
rebuilding long contexts motivate allocating limited KV protection resources
to reduce workflow completion time (WCT).

% One full-width figure with three subfigures; requires subcaption.
% Each landscape panel is 2.24 x 1.85 in at the EuroSys insertion size.
% Plot labels are 10.25 pt at this insertion size; do not shrink this row.
\begin{figure*}[t]
  \centering
  \captionsetup[subfigure]{font=normalsize,labelfont=normalfont,
    textfont=normalfont,justification=centering,singlelinecheck=false}
  \begin{subfigure}[t]{0.32\textwidth}
    \vspace{0pt}
    \includegraphics[width=\linewidth]{figures/motivation-row/motivation-branch-gap.pdf}
    \caption{Branch-duration gaps}
    \label{fig:motivation-branch-gap}
  \end{subfigure}\hfill
  \begin{subfigure}[t]{0.32\textwidth}
    \vspace{0pt}
    \includegraphics[width=\linewidth,height=1.85in,keepaspectratio]{figures/motivation-row/motivation-workflow.drawio.pdf}
    \caption{Four-branch workflow}
    \label{fig:motivation-workflow}
  \end{subfigure}\hfill
  \begin{subfigure}[t]{0.32\textwidth}
    \vspace{0pt}
    \includegraphics[width=\linewidth]{figures/motivation-row/motivation-recovery-exposure.pdf}
    \caption{Recovery cost vs. WCT loss}
    \label{fig:motivation-recovery-exposure}
  \end{subfigure}
  \caption{Motivation for workflow-aware KV protection.
    (a) Each of 121 CC/Weka workflows contributes its largest
    branch-duration gap across strictly overlapping groups with inferred
    fan-in (median: 138\,s; P90: 1,173\,s). Inference uses a Main request
    after all branches finish; Main activity during branches is allowed.
    (b) Four parallel branches with a 91.2\,s duration gap, drawn from
    the same retained groups but not the workflow maxima in (a).
    Time zero is the first branch launch; widths encode durations and
    the right-hand column reports seconds. Dashed tails show slack to
    the 151.6\,s join, determined by A2. Main boxes indicate inferred
    dependencies only. (c) Four independent branch-failure counterfactuals
    at 29.5\,s after the fork. Hatched time is absorbed by slack;
    solid time delays WCT. Costs use Qwen3-14B on A100-SXM with BF16,
    batch one, and TP one, including an assumed 64-token tail. The example illustrates recovery exposure, not
    population-average failure impact.}
  \label{fig:motivation}
  \Description{The panels show workflow-maximum duration gaps, a time-scaled
    four-branch fan-in example with a 91.2-second gap, and estimated
    rebuilding cost decomposed into slack-absorbed time and WCT loss.
    A2 determines the assumed join and exposes its full 32.2-second
    rebuilding cost. A1 absorbs its entire cost; A3 and A4 absorb part
    of theirs, leaving WCT penalties of 17.6 and 17.0 seconds.}
\end{figure*}


\subsection{Workflows, Failures, and KV Recovery}
\label{sec:motivation-background}

\paragraph{Multi-agent workflows.}
We represent a workflow as a directed acyclic graph $G=(V,E)$, whose nodes
are LLM requests and whose edges are data dependencies. An agent branch may
include multiple requests and tool calls; a successor runs only after its
dependencies are satisfied. WCT spans the release of the first request to
the completion of the last. For a given execution schedule, a request's
\emph{slack} is the maximum delay to its completion that does not postpone
workflow completion. Requests on a critical path have zero slack.

\paragraph{Serving and failure model.}
We consider a cluster of GPU workers with block-managed KV that can be
replicated to other GPUs or a host DRAM pool, consistent with KV-centric
architectures such as Mooncake~\cite{qin2025mooncake}. Our failure model is
a fail-stop failure of one GPU or worker: its local KV becomes unavailable,
while the control plane, backup DRAM, and other workers survive. Surviving
workers with the same model already loaded take over affected requests.
Whole-host failures, network partitions, and silent corruption are outside
this failure model.

\paragraph{Recovery paths and resource costs.}
Without a usable copy, the system must rerun prefill to rebuild lost KV.
A DRAM copy avoids reconstructing the saved prefix but must be transferred
to a GPU before use; concurrent recoveries share transfer bandwidth.
A copy already resident in the actual takeover GPU's HBM can be used
directly, at the cost of HBM capacity that could serve foreground requests
and replication traffic before failure. Backup placement therefore trades
capacity and normal-service overhead against recovery speed.

Prior systems reduce failure costs through KV streaming in
D\'{e}j\`{a}Vu~\cite{strati2024dejavu}, erasure coding in
GhostServe~\cite{jayakody2026ghostserve}, and load-aware checkpoint and
recovery coordination in LUMEN~\cite{cao2026lumen}. Multi-agent
applications add an allocation problem: saving the same recovery time for
different requests can have different value to the user. Neither saved KV
bytes nor avoided computation alone captures that difference.

\subsection{Recovery Cost Is Not Workflow Delay}
\label{sec:motivation-workflow}

A join waits for its last required branch. Earlier branches can tolerate
some delay, whereas delaying the branch that determines the join time
postpones downstream work. To characterize branch timing differences, we
analyze the public CC/Weka coding-agent
traces~\cite{semianalysis2026cctraces}. From 393 sessions, we first extract
strictly overlapping subagent groups: the latest branch start precedes
the earliest branch finish. We retain groups whose branches all complete
and are followed by a Main request after the last branch finishes.
Main may also issue requests while branches execute. Because the trace
does not record result dependencies, this criterion infers fan-in from
timing; it does not verify consumption of every branch's output.

The filter retains 286 groups across 121 workflows. For each retained
group, we compute the longest minus the shortest branch duration; each
workflow contributes its maximum such gap.
Figure~\ref{fig:motivation}(\subref{fig:motivation-branch-gap}) shows a median of
138\,s and a 90th percentile of 1,173\,s. These gaps establish substantial
branch-duration heterogeneity. With unequal branch start times, however,
a duration gap is neither join slack nor a direct measure of fault tolerance.

Figure~\ref{fig:motivation}(\subref{fig:motivation-workflow}) shows a
four-branch group from the retained fan-ins. A1--A4 last 60.370, 151.612,
132.409, and 129.429\,s, respectively, giving a duration gap of 91.242\,s.
This group is not its workflow's maximum-gap group, so it does not expand
an individual sample in Figure~\ref{fig:motivation}(\subref{fig:motivation-branch-gap}).
Bars aggregate each branch's requests and tool calls while preserving
observed starts and durations. A2 finishes last at 151.6\,s after the
fork, followed by a Main request at 157.4\,s. Assuming that the successor
depends on all four branches, each branch's slack is its time margin to
A2's completion. Among groups with four or five branches and a gap near
100\,s, we select one with a longer critical-branch prefix to illustrate
recovery exposure, not population-average failure impact.

Figure~\ref{fig:motivation}(\subref{fig:motivation-recovery-exposure}) shows
four independent counterfactuals at 29.5\,s after the fork, each losing
one branch's history KV. The instant lies in an interval with active LLM
requests on all four branches, chosen near the midpoint of their common
execution overlap. Trace hash-block counts provide history lengths
$M_i$ of 100,864--113,664 tokens, to which we add an assumed known
64-token tail. We query a Qwen3-14B static configuration at each exact
length: $T_i$ estimates
empty-cache reconstruction of $M_i+64$ tokens on A100-SXM, with BF16,
batch one, and TP one. The tail state is a counterfactual assumption,
not a trace observation. Slack $\sigma_i$ comes from trace completion
times. Holding other branches, downstream execution, and queueing
unchanged, the WCT penalty is $\max(0,T_i-\sigma_i)$.
This cost substitution does not constitute physical failure replay on
the original serving model.

A1--A4 cost 35.2, 32.2, 36.7, and 38.9\,s to rebuild, but expose
0, 32.2, 17.6, and 17.0\,s to WCT, respectively. A1's 91.2\,s slack
absorbs its entire cost. A2 has the lowest rebuilding cost yet incurs
the largest WCT penalty because it has zero slack. A3 and A4 absorb
only part of their delay with 19.1 and 21.9\,s of slack. Thus, avoiding
the most recomputation need not yield the largest WCT improvement.
KV protection must account for delay tolerance within a workflow;
online estimation and resource contention require the mechanisms and
evaluation developed later.

\subsection{Long Contexts and Recovery Tiers}
\label{sec:motivation-long-context}

% Independent single-column Figure 2; native width 3.34 inches.
% CPU-only static estimates; see reuse-aiconfigurator-max-context-h100/.
\begin{figure}[t]
  \centering
  \includegraphics[width=\linewidth]{figures/motivation-row/motivation-recompute.pdf}
  \caption{Recovery estimates up to each model's context limit.
    Color denotes path; line style and marker denote model.
    Qwen3-14B, Llama-3.1-70B, Qwen3-235B-A22B-Instruct-2507, and DeepSeek-V3.2
    use 1, 4, 16, and 32 H100 GPUs, respectively. The latter two use
    MoE; DeepSeek (671B) also uses DSA. BF16, batch one.
    Total context includes the known 64-token tail; endpoints are
    128K (YaRN), 128K, 1,010,000, and 160K tokens, respectively.
    Axis suffixes: K~=~1,024; M~=~$10^6$. Recompute uses sequential 8K chunks.
    DRAM restores full per-rank state concurrently at an assumed
    64\,GB/s per GPU before compute, without overlap.
    All curves are static estimates. Beyond 256K, Qwen 2507 uses a
    full-attention cost projection without DCA or sparse-kernel acceleration.
    DSA estimates use analytical prefix correction and table extrapolation
    beyond 32K. Context limits do not imply validated recovery or accuracy.
    Detection, queueing, loading, pre-failure copying, and software
    takeover overheads are excluded.}
  \label{fig:motivation-recompute}
  \Description{One plot compares twelve estimated recovery curves on
    logarithmic context and latency axes. Blue, orange, and red encode
    HBM reuse, DRAM reuse, and Recompute. Circles, squares, triangles,
    and diamonds encode Qwen3-14B, Llama-3.1-70B, Qwen3-235B-A22B-Instruct-2507,
    and DeepSeek-V3.2, respectively. The latter two are MoE models;
    DeepSeek additionally uses sparse attention. Curves stop at each
    model's configured context limit. At 1,010,000 total tokens,
    Qwen 2507 costs an estimated 70 milliseconds for HBM reuse,
    830 milliseconds for DRAM reuse, and 220 seconds for chunked
    full-attention recomputation. GPU allocations differ across models.}
\end{figure}


Repeated agent interactions accumulate code, tool outputs, and other
agents' intermediate results. Backups avoid repeating that computation,
but their location still affects recovery latency.
Figure~\ref{fig:motivation-recompute} compares three initial states:
all historical KV and required auxiliary state reside in the takeover
GPU group's HBM, the same state is in host-local DRAM, or no copy is
available. All paths end after the same known 64-token tail's forward
pass, with history and tail state ready for decoding; weights are
already loaded. HBM reuse therefore still incurs tail computation.

We use AIConfigurator~\cite{nvidia2026aiconfigurator} 0.9.0 HYBRID estimates
for dense Qwen3-14B and Llama-3.1-70B, the Instruct-2507 version of
Qwen3-235B-A22B (235B total, 22B active parameters) with MoE, and DeepSeek-V3.2
(671B total, 37B active) with MoE and DeepSeek Sparse Attention (DSA).
MoE selects feed-forward experts. We model Qwen with dense GQA and
DSA with its historical-token indexer. We use
H100-SXM 80\,GiB, vLLM 0.19.0, BF16 weights and KV, and batch one.
TP is 1, 4, 16, and 32, respectively; the MoE models use EP 16 and 32,
and all use PP one. GPU allocations differ, so we compare recovery
paths within each model configuration.

The horizontal axis is total context $C=M+64$, including the known tail.
Each curve reaches its model-specific limit: 131,072 tokens with YaRN
for Qwen3-14B~\cite{qwen2025technical}, 131,072 for
Llama-3.1-70B~\cite{meta2024llama31card}, 163,840 in the official
DeepSeek-V3.2 configuration~\cite{deepseek2025v32config}, and 1,010,000
in Qwen 2507's extended configuration~\cite{qwen2025instruct2507}.
Qwen 2507 natively supports 256K; beyond that length, we project
full-attention costs without modeling Dual Chunk Attention (DCA)
position handling or sparse-kernel acceleration. Those points do not
estimate the performance of the official sparse 1M serving path.

HBM reuse computes only the 64-token tail. Recompute starts empty and
processes sequential chunks of at most 8,192 tokens, reusing KV rebuilt
by earlier chunks; vocabulary projection is excluded
from all paths because it does not rebuild KV. A common chunk size bounds modeled
activation memory, rather than interpreting an unchunked-prefill
capacity warning as a model's context limit. Inter-chunk scheduling
overhead is excluded. DRAM first restores each rank's full historical
state, then executes the same tail. GQA transfer bytes account for KV
sharding and head replication. DSA includes replicated MLA latents,
RoPE keys, and native FP8 indexer keys with FP32 scales within the
otherwise BF16 configuration. We assume concurrent restoration at
64\,GB/s per rank, with no shared host bottleneck or transfer--compute
overlap. This bandwidth is a system-model assumption, not a measurement.

At 128K total context, HBM reuse costs an estimated 22--56\,ms,
DRAM reuse 155--357\,ms, and chunked Recompute 14.5--21.0\,s.
At 1,010,000 tokens, Qwen 2507's full-attention cost projection gives
70\,ms, 830\,ms, and 220\,s, respectively. For Qwen3-14B at 128K,
a DRAM copy avoids most historical recomputation; full HBM residency
additionally removes about 335\,ms of modeled transfer waiting.
Thus, backups avoid computation, while HBM's incremental value over
DRAM comes from reducing transfer waiting. Parameter counts, state
representation, and parallelism jointly affect the cost.

HYBRID combines kernel tables and performance models; these are not
physical recovery measurements. DSA module tables end at 32K total
tokens, so longer shapes involve extrapolation; cached-prefix queries
use analytical prefix correction. All 25 model/context combinations
pass the static memory check, with a maximum estimate of 77.8\,GiB
per GPU, without validating deployment, long-context accuracy, or
native sparse recovery. Full-history HBM residency is a reference,
not AegisKV's limited hot-block placement. These estimates omit
contention and cannot establish WCT gains under limited backup capacity.
Figure~\ref{fig:motivation} further shows why recovery savings must be
interpreted in the context of branch delay tolerance.

\subsection{Problem and Design Goals}
\label{sec:motivation-goals}

AegisKV asks: \emph{with limited backup capacity and replication and
recovery bandwidth, which KV from which requests should receive what
protection to reduce the workflow-level impact of failures?}
Dependencies differentiate protection benefits, long contexts increase
state-loss costs, and replication competes with foreground serving.

This problem imposes five goals and constraints.
\textbf{(R1)} Optimize recovery for WCT, accounting for dependency-induced
delay tolerance.
\textbf{(R2)} Allocate limited backup capacity across requests and their
KV to reduce exposed recovery waiting.
\textbf{(R3)} Support concurrent recovery and transfer--compute overlap
without introducing request-level serialization for protection priority.
\textbf{(R4)} Bound normal-service replication overhead and subordinate
backup capacity and bandwidth use to foreground needs.
\textbf{(R5)} Preserve model computation and KV recovery correctness using
existing serving resources.
The next section develops the evidence for mechanism choices and the
corresponding placement and recovery design.
